\documentclass[11pt]{article}
\usepackage[a4paper,margin=1in]{geometry}
\usepackage{fontspec}
\usepackage[boldfont=false]{xeCJK}
\setCJKmainfont{FandolSong-Regular.otf}
\usepackage{amsmath}
\usepackage{amssymb}
\usepackage{array}
\usepackage{booktabs}
\usepackage{graphicx}
\usepackage{adjustbox}
\usepackage{enumitem}
\setlist[itemize]{topsep=2pt,partopsep=0pt,parsep=0pt,itemsep=2pt,leftmargin=1.4em}
\setlist[enumerate]{topsep=2pt,partopsep=0pt,parsep=0pt,itemsep=2pt,leftmargin=1.6em}
\usepackage{parskip}
\usepackage{fvextra}
\fvset{breaklines=true,breakanywhere=true,fontsize=\small}
\setlength{\parindent}{0pt}

\begin{document}

\begin{center}
  {\LARGE\bfseries Business and its environment}\\[0.35em]
  {\large A Level Business}
\end{center}
\vspace{0.6em}

\section*{Why businesses exist}

People have \textbf{needs} 需要 (things you must have, like food and shelter) and \textbf{wants} 欲望 (things you would like, like a phone or a holiday). A business exists to meet these needs and wants by making \textbf{goods} 商品 and \textbf{services} 服务.

To do this, a business uses \textbf{resources} 资源. We sort resources into four \textbf{factors of production} 生产要素:

\begin{itemize}
\item \textbf{land} 土地 — natural things, such as soil, water and minerals.

\item \textbf{labour} 劳动力 — the work done by people.

\item \textbf{capital} 资本 — the machines, tools, buildings and money used to produce.

\item \textbf{enterprise} 创业精神 — the skill of bringing the other three together and taking the \textbf{risk} 风险 to start a business.

\end{itemize}

Resources are \textbf{scarce} 稀缺 — there are not enough to meet everyone's wants. So every business must choose how to use them well.

\section*{Adding value}

\textbf{Adding value} 增值 means the selling price of a product is higher than the cost of the \textbf{inputs} 投入 (the bought-in materials) used to make it.

\[
\text{value added} = \text{selling price} - \text{cost of bought-in materials}
\]

A business can add value by \textbf{branding} 品牌, better quality, good design, fast service, or convenience. More value added can mean more \textbf{profit} 利润, but only if costs stay under control.

\section*{Entrepreneurs and intrapreneurs}

An \textbf{entrepreneur} 企业家 is a person who takes the risk to set up and run a business. They have the idea, organise the factors of production, and make the decisions. If the business fails, they can lose their money.

An \textbf{intrapreneur} 内部创业者 is an employee inside a larger business who acts like an entrepreneur. They build new ideas, products or ways of working, using the company's resources.

\subsection*{Characteristics of successful entrepreneurs}

Successful entrepreneurs are usually:

\begin{itemize}
\item \textbf{innovative} 有创新精神的 — full of new ideas.

\item willing to take risks, but in a careful way.

\item hard-working and \textbf{determined} 有决心的 — they do not give up easily.

\item \textbf{self-confident} 自信的 and good at making decisions.

\item good leaders who can plan ahead and manage money.

\end{itemize}

\subsection*{Why enterprise matters to a country}

When people start new businesses (\textbf{business start-ups} 初创企业), the whole country can gain:

\begin{itemize}
\item new jobs are created, which lowers \textbf{unemployment} 失业.

\item more goods and services are made, which raises \textbf{output} 产出 (and the country's \textbf{GDP} 国内生产总值).

\item more \textbf{competition} 竞争, giving customers lower prices and more choice.

\item some firms grow and \textbf{export} 出口, bringing money into the country.

\item \textbf{economic growth} 经济增长 and a higher standard of living.

\end{itemize}

\section*{The business plan}

A \textbf{business plan} 商业计划书 is a written document that describes the business idea and how it will work. It usually contains:

\begin{center}
\adjustbox{max width=\linewidth}{%
\begin{tabular}{ll}
\toprule
\textbf{Part of the plan} & \textbf{What it covers} \\
\midrule
The idea & the product, and what makes it different \\
The market & target customers, market size, competitors \\
Marketing & price, promotion and place \\
Operations & how the product is made or delivered \\
Finance & start-up costs, a \textbf{cash-flow forecast} 现金流量预测, expected profit \\
People & the owners and key staff \\
\bottomrule
\end{tabular}}
\end{center}

A business plan is useful because it helps the owner:

\begin{itemize}
\item get \textbf{finance} 融资 from a bank or an \textbf{investor} 投资者.

\item set clear goals and lower the risk of failure.

\item guide decisions as the business grows.

\end{itemize}

\section*{Sectors of business activity}

Every business works in one of four sectors:

\begin{center}
\adjustbox{max width=\linewidth}{%
\begin{tabular}{lll}
\toprule
\textbf{Sector} & \textbf{What it does} & \textbf{Example} \\
\midrule
\textbf{primary sector} 第一产业 & takes raw materials from nature & farming, mining, fishing \\
\textbf{secondary sector} 第二产业 & makes (manufactures) goods & car factory, baker \\
\textbf{tertiary sector} 第三产业 & provides services & shops, banks, transport \\
\textbf{quaternary sector} 第四产业 & knowledge and information services & research, IT, consulting \\
\bottomrule
\end{tabular}}
\end{center}

Here \textbf{raw materials} 原材料 are the natural inputs a business starts with, and to \textbf{manufacture} 制造 means to make goods in large amounts. As a country develops, it usually moves from mostly primary work towards more tertiary and quaternary work.

\section*{Private and public sectors}

\begin{itemize}
\item The \textbf{private sector} 私营部门 is made up of businesses owned by private people. Their main aim is usually profit.

\item The \textbf{public sector} 公共部门 is made up of organisations owned by the government, such as state schools and public hospitals. Their aim is usually to provide a service, not to make profit.

\end{itemize}

\section*{Legal structures of business}

A business must choose a legal form. The main ones are below.

\begin{itemize}
\item A \textbf{sole trader} 个体经营者 is owned by one person. It is easy to set up and the owner keeps all the profit, but raising money is hard.

\item A \textbf{partnership} 合伙企业 is owned by two or more partners who share the work and profit. They often sign a \textbf{deed of partnership} 合伙协议 that sets out each partner's share.

\item A \textbf{private limited company} 私人有限公司 (Ltd) sells \textbf{shares} 股份 to a small group, such as family and friends. Its owners are \textbf{shareholders} 股东. Its shares are not sold to the public.

\item A \textbf{public limited company} 上市公司 (plc) sells its shares to the public on the \textbf{stock exchange} 证券交易所. It can raise large amounts of money, but it must follow more rules and faces the risk of a \textbf{takeover} 收购.

\item A \textbf{franchise} 特许经营 lets one firm use another firm's brand and business model. The \textbf{franchisor} 特许方 owns the brand; the \textbf{franchisee} 加盟方 pays a fee and a \textbf{royalty} 提成 to run a branch.

\item A \textbf{joint venture} 合资企业 is when two businesses join to run one project and share the costs and risks.

\item A \textbf{social enterprise} 社会企业 trades like a normal business but aims to do social or \textbf{environmental} 环境 good, and reinvests most of its profit.

\end{itemize}

\section*{Limited and unlimited liability}

\textbf{Liability} 责任 means who must pay the business's \textbf{debts} 债务 if it cannot pay them itself.

\begin{itemize}
\item \textbf{limited liability} 有限责任 — the owners can only lose the money they put in. Their personal things, like a house or car, are safe. Companies have limited liability.

\item \textbf{unlimited liability} 无限责任 — the owners are personally responsible for all the debts. They could lose personal things. Sole traders and most partnerships have unlimited liability.

\end{itemize}

A company is a \textbf{separate legal identity} 独立法人 — it can own things, owe money, and be taken to court in its own name, apart from its owners.

\section*{The size of a business}

There is no single best way to measure how big a business is. Common measures are:

\begin{itemize}
\item the number of \textbf{employees} 员工.

\item \textbf{revenue} 营业收入 (also called sales turnover) — the value of what it sells.

\item \textbf{capital employed} 所用资本 — the total money invested in the business.

\item \textbf{market share} 市场份额 — the firm's sales as a percentage of the whole market.

\item \textbf{market capitalisation} 市值 — the total value of a company's shares (for a plc).

\end{itemize}

Different measures can give different answers, so it is best to use more than one.

\subsection*{Why small businesses matter}

Small businesses create jobs, serve local \textbf{communities} 社区, give personal service, react quickly to change, bring new ideas, and supply parts and services to large firms.

\section*{Business growth}

A business can grow in two ways.

\textbf{Internal growth} 内部增长 (organic growth) — the business grows by itself: selling more, opening new branches, or adding new products. It is slower but lower risk.

\textbf{External growth} 外部增长 — the business joins with another firm. This is called \textbf{integration} 整合, and it happens by:

\begin{itemize}
\item a \textbf{merger} 合并 — two firms agree to join together as one.

\item a \textbf{takeover} — one firm buys control of another.

\end{itemize}

There are three types of integration:

\begin{itemize}
\item \textbf{horizontal integration} 横向整合 — joining a firm at the same stage in the same industry (e.g. two car makers).

\item \textbf{vertical integration} 纵向整合 — joining a firm at a different stage in the same industry. \textbf{Forward integration} 前向整合 is towards the customer (e.g. buying a shop); \textbf{backward integration} 后向整合 is towards the supplier (e.g. buying a parts maker).

\item \textbf{conglomerate integration} 混合整合 — joining a firm in a completely different industry. This spreads risk, which is called \textbf{diversification} 多元化.

\end{itemize}

Firms grow to lower their costs, win a bigger market share and more \textbf{market power} 市场支配力, raise profit, and spread risk. But growth also brings problems: higher average costs, harder management and communication, different company \textbf{cultures} 企业文化 that may clash, and a need for a lot of finance.

\section*{Economies and diseconomies of scale}

\textbf{Economies of scale} 规模经济 are the cost savings a business gets as it grows larger. The \textbf{average cost} 平均成本 of each unit falls as output rises. The main types are:

\begin{center}
\adjustbox{max width=\linewidth}{%
\begin{tabular}{ll}
\toprule
\textbf{Type} & \textbf{How it lowers cost} \\
\midrule
purchasing & buying in \textbf{bulk} 大批量 earns discounts \\
technical & large machines are more efficient \\
financial & large firms borrow at lower \textbf{interest} 利息 \\
managerial & the firm can employ expert specialist managers \\
marketing & advertising cost is spread over more units \\
\bottomrule
\end{tabular}}
\end{center}

\textbf{Diseconomies of scale} 规模不经济 happen when a firm grows too big and the average cost starts to rise again. The usual causes are:

\begin{itemize}
\item \textbf{communication} 沟通 becomes slow and unclear.

\item the \textbf{coordination} 协调 of many departments is hard.

\item workers may feel unimportant and lose \textbf{motivation} 激励.

\end{itemize}

\section*{Aims, objectives and mission}

A \textbf{mission statement} 使命宣言 is a short sentence that says why the business exists and what it wants to be. It guides everyone in the firm.

From the mission come \textbf{aims} 目的 (broad, long-term goals), and then \textbf{objectives} 目标 (the specific targets that help reach the aims).

Good objectives are \textbf{SMART} — the first letters of:

\begin{itemize}
\item \textbf{Specific} 具体的 — clear and exact.

\item \textbf{Measurable} 可衡量的 — you can check it with numbers.

\item \textbf{Achievable} 可实现的 — it is possible to reach.

\item \textbf{Realistic} 现实的 — sensible for the resources you have.

\item \textbf{Time-bound} 有时限的 — it has a deadline.

\end{itemize}

A \textbf{strategy} 战略 is the long-term plan to meet objectives. \textbf{Tactics} 策略 are the short-term actions that carry out the strategy. Objectives are shaped by the size and age of the firm, who owns it, the company culture, and conditions in the market.

\section*{Types of business objective}

\begin{center}
\adjustbox{max width=\linewidth}{%
\begin{tabular}{ll}
\toprule
\textbf{Objective} & \textbf{What it means} \\
\midrule
\textbf{profit maximisation} 利润最大化 & making the largest possible profit \\
growth & getting bigger — more sales, more branches \\
\textbf{survival} 生存 & staying in business, vital for new firms and in hard times \\
market share & winning a larger share of the market \\
social and \textbf{ethical} 道德 objectives & doing what is morally right, such as fair pay \\
\bottomrule
\end{tabular}}
\end{center}

\textbf{Corporate social responsibility} 企业社会责任 (CSR) is the idea that a business should act in the best interests of society and the environment, not just make profit. Examples are treating workers fairly, cutting \textbf{pollution} 污染, and helping the local community.

A firm's objectives can change over time. A new firm may aim to survive; later it may aim to grow; a large, safe firm may focus on profit or CSR.

\section*{Stakeholders}

A \textbf{stakeholder} 利益相关者 is any person or group with an interest in a business and how well it does.

\textbf{Internal stakeholders} 内部利益相关者 are inside the business:

\begin{itemize}
\item \textbf{owners} and shareholders — want profit and a good return.

\item \textbf{managers} 管理者 — want growth, job security and bonuses.

\item employees — want fair pay, safe work and job security.

\end{itemize}

\textbf{External stakeholders} 外部利益相关者 are outside the business:

\begin{center}
\adjustbox{max width=\linewidth}{%
\begin{tabular}{ll}
\toprule
\textbf{Stakeholder} & \textbf{Main objective} \\
\midrule
\textbf{customers} 顾客 & good quality at a fair price \\
\textbf{suppliers} 供应商 & regular orders, paid on time \\
\textbf{lenders} 贷款方 (banks) & to be repaid with interest \\
the government & tax, jobs, and that the firm obeys the law \\
local community & jobs, but little pollution or noise \\
\textbf{pressure groups} 压力团体 & to protect a cause, such as the environment \\
\bottomrule
\end{tabular}}
\end{center}

\subsection*{Why stakeholder objectives conflict}

Stakeholders often want different things, so their objectives can \textbf{conflict} 冲突 (clash). For example:

\begin{itemize}
\item owners want higher profit, but employees want higher pay — and higher pay lowers profit.

\item customers want low prices, but owners want a high profit margin.

\item the local community wants less pollution, but cutting pollution can raise costs.

\end{itemize}

\subsection*{How conflict can be managed}

\begin{itemize}
\item keep good communication with every group.

\item \textbf{compromise} 妥协 — give each group part of what it wants.

\item \textbf{negotiation} 谈判 — talk until the groups reach an agreement.

\item decide which stakeholders matter most for each decision, and act on that.

\end{itemize}


\end{document}
